\chapter{RDW (Red Cell Distribution Width)}

\section*{What Is This Test?}
The red cell distribution width (RDW) is a quantitative measure of anisocytosis --- the variation in size of circulating red blood cells. It is reported as either the coefficient of variation (RDW-CV), expressed as a percentage, or the standard deviation (RDW-SD), expressed in femtoliters. A normal RDW indicates a relatively uniform population of red blood cells, while an elevated RDW reflects a heterogeneous mixture of large and small cells. The RDW is automatically calculated by hematology analyzers from the red blood cell volume distribution histogram and is reported as part of the standard complete blood count. The RDW is particularly useful in the differential diagnosis of anemia when paired with the mean corpuscular volume (MCV). For example, elevated RDW is an early and sensitive marker of iron deficiency anemia, often preceding the fall in hemoglobin and MCV. In contrast, thalassemia trait typically presents with microcytosis (low MCV) but a normal RDW because the genetic defect produces uniformly small red cells. This MCV/RDW classification system (Bessman classification) provides a systematic approach to the differential diagnosis of anemia. Beyond hematology, an elevated RDW has emerged as a powerful independent predictor of adverse outcomes in cardiovascular disease, heart failure, critical illness, and overall mortality.

\section*{Alternative Names}
RDW; Red Cell Distribution Width; RDW-CV; RDW-SD; Anisocytosis Index; Erythrocyte Distribution Width; Red Blood Cell Distribution Width

\section*{Principle}
The RDW is derived from the red blood cell volume histogram generated by automated hematology analyzers using either the electrical impedance (Coulter) principle or optical laser light scatter. As each erythrocyte passes through the detection aperture or laser beam, its volume is measured. The analyzer records thousands to tens of thousands of individual cell volume measurements and constructs a frequency distribution histogram of red blood cell volumes. From this histogram, the statistical distribution parameters are calculated. RDW-CV is calculated as: RDW-CV (\%) = (Standard deviation of RBC volume / MCV) $\times$ 100. A normal RDW-CV is approximately 11.5--14.5\%. RDW-SD is the actual width (in fL) of the erythrocyte volume distribution curve measured at 20\% of the curve height. Normal RDW-SD values are approximately 38--47 fL but are analyzer-specific. The CV computation is more sensitive to the MCV value (since it is the denominator), whereas the SD value represents a direct linear measurement of the volume distribution width. Both are reported by most modern hematology analyzers, including Beckman Coulter DxH, Sysmex XN, and Siemens ADVIA platforms.

\section*{History}
The concept of measuring red blood cell size variation dates back to the visual estimation of anisocytosis on peripheral blood smears by early hematologists. However, a quantitative numerical index was not possible until automated analyzers capable of measuring individual cell volumes became available. The Coulter Counter Model S (introduced in 1968) was the first instrument to generate red cell volume distribution histograms. The term ``red cell distribution width'' and the abbreviation RDW were popularized by Dr. J. David Bessman and colleagues in a landmark 1983 paper in the American Journal of Clinical Pathology, which demonstrated that combining MCV and RDW could classify anemias with greater specificity than MCV alone. This Bessman classification became the standard approach to initial anemia workup and remains widely taught in medical education today. The recognition that RDW elevation is an independent predictor of mortality across diverse patient populations emerged in the early 2000s and has made the RDW a parameter of interest well beyond hematology.

\section*{Why This Test Is Ordered}
\begin{itemize}
  \item Early detection of iron deficiency anemia --- RDW is often elevated before hemoglobin falls or MCV decreases, making it a sensitive screening marker for iron-deficient erythropoiesis
  \item Distinction between iron deficiency anemia (elevated RDW, low MCV) and thalassemia trait (normal RDW, low MCV) in patients with microcytic anemia
  \item Classification of anemia using the Bessman MCV/RDW system: microcytic homogeneous (thalassemia trait), microcytic heterogeneous (iron deficiency), normocytic homogeneous (anemia of chronic disease), normocytic heterogeneous (early iron/B12/folate deficiency, mixed deficiency), macrocytic homogeneous (aplastic anemia), macrocytic heterogeneous (megaloblastic anemia)
  \item Identification of dimorphic anemia --- mixed iron and B12/folate deficiency produces markedly elevated RDW
  \item Monitoring response to hematinic therapy --- successful iron, B12, or folate replacement transiently increases RDW as new (normal-sized) cells mix with existing abnormal cells, then RDW normalizes
  \item Prognostic assessment --- elevated RDW predicts increased all-cause mortality, cardiovascular events, and adverse outcomes in heart failure, coronary artery disease, stroke, critical illness, and community-dwelling older adults
\end{itemize}

\section*{Is This a Primary or Secondary Test}
The RDW is a primary calculated parameter reported with the CBC and is essential in the classification and differential diagnosis of anemia.

\section*{How This Test Is Performed}
The RDW is calculated automatically from a venous blood sample collected in a lavender-top (EDTA) tube and processed on an automated hematology analyzer. No separate sample or additional testing is required beyond the standard CBC. A healthcare professional draws approximately 2--5 mL of venous blood, typically from the antecubital fossa. The sample is gently inverted to mix with EDTA and placed on the analyzer. As red blood cells pass through the detection system, the volume of each cell is measured, and the volume distribution histogram and RDW are automatically computed. Venipuncture takes less than 5 minutes; automated analysis is completed within 1--2 minutes.

\section*{What Preparation Is Needed}
No special preparation is required. Fasting is not necessary for RDW measurement. Patients should avoid vigorous exercise immediately before blood collection. Medications and supplements, particularly iron, vitamin B12, folate, and recent blood transfusion, should be disclosed to the healthcare provider, as they can affect the RDW.

\section*{Contraindications and Precautions}
No absolute contraindications. Factors affecting RDW accuracy include: very small sample volume (inadequate for accurate histogram generation), hemolyzed or clotted specimens (invalid), delayed sample processing (>24 hours, causing red cell swelling that alters the volume distribution), cold agglutinins (causes red cell clumping, distorting the volume histogram and spurious RDW elevation), recent blood transfusion (produces a bimodal distribution with elevated RDW), and marked leukocytosis (>100,000/$\mu$L may distort the red cell histogram). Elevated RDW in the absence of hematologic disease (anemia) may reflect underlying inflammatory, nutritional, or chronic disease states. EDTA is the standard anticoagulant. Samples should be analyzed within 4 hours at room temperature or 24 hours refrigerated at 2--8\textdegree{}C.

\section*{Risks}
The risks are those of routine venipuncture: mild pain or stinging at the needle site, bruising or small hematoma formation, lightheadedness or vasovagal syncope, and rarely, infection or nerve injury.

\section*{What Happens After the Test}
RDW results are available within 30--60 minutes as part of the CBC report. The healthcare provider interprets the RDW in combination with the MCV and hemoglobin. An elevated RDW with microcytosis prompts iron studies. A normal RDW with microcytosis prompts hemoglobin electrophoresis for thalassemia trait. An elevated RDW with macrocytosis prompts B12 and folate levels. The Bessman classification table (MCV vs. RDW) is used to narrow the diagnostic possibilities. Serial RDW measurements are monitored during hematinic therapy to track the replacement of abnormal cells with normal ones. In non-anemic patients with unexplained RDW elevation, evaluation for occult inflammation, chronic kidney disease, liver disease, or nutritional deficiency may be pursued.

\section*{Understanding Results}

\subsection*{Below Normal (RDW-CV <11.5\%)}
\begin{itemize}
  \item \textbf{Thalassemia trait:} Uniformly microcytic red blood cells with normal RDW and low MCV. The genetic defect in globin chain synthesis produces a consistent, homogeneous microcytosis. This pattern, combined with normal iron studies and elevated HbA2 on hemoglobin electrophoresis, is diagnostic of beta-thalassemia trait.
  \item \textbf{Anemia of chronic disease (early):} Mild normocytic anemia with uniform cell population and normal RDW.
  \item \textbf{Hereditary spherocytosis:} Uniform spherical cells with normal RDW, elevated MCHC, and spherocytes on peripheral smear.
  \item \textbf{No pathological significance:} A low or normal RDW in the absence of hematologic abnormalities is a normal finding and does not indicate disease.
\end{itemize}

\subsection*{Above Normal (RDW-CV >14.5\%)}
\begin{itemize}
  \item \textbf{Iron deficiency anemia:} RDW is often the earliest laboratory parameter to become abnormal in iron deficiency, rising before the hemoglobin or MCV decline. The anisocytosis reflects a heterogeneous mixture of normal and microcytic hypochromic cells as iron stores are depleted. RDW can exceed 20\% in severe iron deficiency. Serum ferritin, transferrin saturation, and TIBC confirm the diagnosis.
  \item \textbf{Megaloblastic anemia (B12 or folate deficiency):} Macro-ovalocytes coexist with normocytes and microcytes, producing a markedly elevated RDW along with elevated MCV. Hypersegmented neutrophils on peripheral smear, low serum B12 or folate, and elevated methylmalonic acid/homocysteine confirm the diagnosis.
  \item \textbf{Mixed deficiency anemia:} Combined iron deficiency (microcytosis) and B12/folate deficiency (macrocytosis) produces a dimorphic blood picture with the highest RDW values seen clinically, often 20--30\%. The peripheral smear shows both microcytic hypochromic cells and macro-ovalocytes.
  \item \textbf{Myelodysplastic syndromes:} Dyserythropoiesis with marked anisocytosis, seen with macrocytosis and cytopenias in other lineages. Bone marrow examination with cytogenetic analysis confirms the diagnosis.
  \item \textbf{Hemolytic anemia:} Reticulocytosis introduces a population of larger, younger cells that mix with mature normocytes, elevating the RDW.
  \item \textbf{Post-treatment response:} Following iron, B12, or folate supplementation, newly produced normal-sized cells mix with the pre-existing abnormal population, transiently increasing RDW before normalization.
  \item \textbf{Post-transfusion:} Transfused donor cells of different age and size mixed with the recipient's native population.
  \item \textbf{Chronic liver disease:} Altered lipid metabolism and splenic sequestration produce anisocytosis.
  \item \textbf{Hemoglobinopathies:} Sickle cell disease, HbC disease, and related disorders.
  \item \textbf{Unexplained elevated RDW (no anemia):} Associated with cardiovascular disease, heart failure, chronic kidney disease, inflammatory conditions (elevated CRP), malnutrition, and aging; predicts increased all-cause mortality.
\end{itemize}

\section*{Normal Results}
\begin{itemize}
  \item \textbf{RDW-CV (adults):} 11.5--14.5\%
  \item \textbf{RDW-SD (adults):} 38--47 fL (analyzer-dependent; check laboratory-specific reference ranges)
  \item \textbf{Neonates:} 14.5--18.5\% (physiological anisocytosis is normal in the first weeks of life)
  \item \textbf{Children (1--12 years):} 11.5--14.5\% (similar to adult range after the neonatal period)
\end{itemize}

\section*{Follow-up Tests}
\begin{itemize}
  \item \textbf{Peripheral blood smear:} Direct visual assessment of anisocytosis, poikilocytosis, microcytosis, macrocytosis, hypochromia, target cells, and red cell inclusions
  \item \textbf{Iron studies (serum iron, ferritin, total iron-binding capacity, transferrin saturation):} Essential for confirming iron deficiency when RDW and MCV are consistent with iron-deficient erythropoiesis
  \item \textbf{Vitamin B12 and serum folate levels:} When RDW is elevated with macrocytosis
  \item \textbf{Methylmalonic acid and serum homocysteine:} More sensitive tests for B12 and folate deficiency when serum levels are borderline
  \item \textbf{Hemoglobin electrophoresis:} To diagnose thalassemia trait (elevated HbA2) or other hemoglobinopathies when microcytosis is present with normal RDW
  \item \textbf{Reticulocyte count:} To assess for reticulocytosis contributing to RDW elevation
  \item \textbf{Serum erythropoietin level:} In unexplained RDW elevation
  \item \textbf{C-reactive protein (CRP) and erythrocyte sedimentation rate (ESR):} To assess for underlying inflammation contributing to elevated RDW
  \item \textbf{Liver function tests and renal function tests:} To evaluate for chronic liver or kidney disease as contributors to RDW elevation
  \item \textbf{Bone marrow aspiration and biopsy:} If myelodysplastic syndrome, aplastic anemia, or other marrow disorder is suspected
\end{itemize}

\section*{References}
\begin{enumerate}
  \item Bessman JD, Gilmer PR Jr, Gardner FH. Improved classification of anemias by MCV and RDW. Am J Clin Pathol. 1983;80(3):322--326.
  \item Tierney LM, Saint S, Drazen JM. CURRENT Medical Diagnosis \& Treatment. McGraw-Hill; 2023.
  \item Wallach J. Interpretation of Diagnostic Tests. 11th ed. Wolters Kluwer; 2022.
  \item Fischbach FT, Dunning MB. A Manual of Laboratory and Diagnostic Tests. 11th ed. Wolters Kluwer; 2023.
  \item Burtis CA, Ashwood ER, Bruns DE. Tietz Textbook of Clinical Chemistry and Molecular Diagnostics. 7th ed. Elsevier; 2023.
  \item Kaushansky K, Lichtman MA, Prchal JT, et al. Williams Hematology. 10th ed. McGraw-Hill; 2021.
  \item Lippi G, Plebani M. Red blood cell distribution width (RDW) and human pathology: one size fits all. Clin Chem Lab Med. 2014;52(9):1247--1249.
  \item Danese E, Lippi G, Montagnana M. Red blood cell distribution width and cardiovascular diseases. J Thorac Dis. 2015;7(10):E402--E411.
\end{enumerate}

\clearpage
